\section{Ten-positive equations by denominator engineering}\label{sec:compact}
This section supplies six complete ordinary integer polynomials, each in ten
strictly positive integer unknowns. Four give useful local operation/degree
tradeoffs; two are retained as dominated controls. All fixed coefficients are
the \emph{whole} indexed BFS943 native tuple, supplied literally in
\texttt{compact/coefficients/}. Its finite reconstruction is
\texttt{rebuild\_coefficients.py}; its largest port has 592,884 bits.
These are not coefficient-free formulas, and fixed ports are not unknowns.

\begin{center}
\begin{tabular}{llrr}
Chart & Weight & Binary operations & Exact degree\\\hline
$Y=q^3s,\ D=XY$ & $V^2$ &77&626\\
$Y=q^3s,\ D=XY$ & $bV$ &78&464\\
$Y=q^3s,\ D=q^4$ & $V^2$ &78&622\\
$Y=q^3s,\ D=q^4$ & $bV$ &79&460\\
$Y=q^2s,\ D=q^3Y$ & $V^2$ &79&518\\
$Y=q^2s,\ D=q^3Y$ & $bV$ &80&384
\end{tabular}
\end{center}
The $78/622$ and $79/518$ points are dominated by $78/464$; they are
uniform-family controls, not frontier claims. The incoming unscaled
$78/630$ and $79/468$ points used $D=Y^2$. The new $77/626$ and $78/464$
points respectively improve both metrics of those matched predecessors.
The $80/384$ point also improves the incoming unshifted $83/436$ point
in both metrics. Older lower-degree constructions remain relevant.
No global optimum or literature-wide priority is claimed.

\paragraph{Attribution and scope.}
The incoming \emph{Certificate Reductions} argument already removed the
artificial scale and recovered a Pell index from a known divisor before
recovering a whole denominator. \emph{Semantic Signs} already proved
unshifted transport and the exclusion of both signed third auxiliary
branches. Pell doubling, product-square extraction, and quotient
normalization are inherited. Here the new choices are $D=XY$, $D=q^4$,
and the partially normalized $Y=q^2s,D=q^3Y$, whose fractional carrier
numerator reuses $Y^2$. The proof below composes explicitly with the
retained ordinary Pell/source theorems. It is not a new end-to-end
proof-assistant certificate and proves neither consistency nor
unsolvability separately.

\subsection{The complete polynomials}
The only independent coordinates are
\[
 J,F,\alpha,i,T,w,s,n,\eta,\zeta\in\Z_{>0}.
\]
The fixed positive integers $\beta,K_0,M_C,M_F$ are the unchanged indexed
tuple, with $\beta\ge31$ odd and $0<M_C<\beta<M_F<2\beta$.
Define polynomial abbreviations
\[
\begin{gathered}
q=\beta J+1,\quad h=q-1,\quad C=q-F-\alpha,\quad X=wq,\\
k=\eta+\zeta,\quad c=kY+\eta,\quad A=Y(X+1)+1,\\
\Delta=A^2-1,\quad H=4A-5,\quad v=XY^2,\\
E=1+v(v+1)k^2,\\
R=[h(q-F)+\alpha](q^2-1)+(M_C+qM_F)J,\\
e=k-R-1,\quad B=(K_0+w)C-F,\\
U=ic^2,\quad a=\Delta U^2E,\quad F_0=1+a,\quad G=EF_0,\\
F_{\rm sq}=1+4aF_0,\quad Q=\Delta(F_{\rm sq}-1),\\
b=cT-R,\quad V=F_{\rm sq}b-c,\quad K=X(3c-2X)+2.
\end{gathered}
\]
Choose $(Y,D)$ from $(q^3s,XY),(q^3s,q^4),(q^2s,q^3Y)$ and choose
$L=V^2$ or $L=bV$. Put $M=hDH$. In the first two charts set
\[
 Z=hD\{H(n-LB)+K\}-H(DB+he).
\]
In the partial chart $Y=q^2s,D=q^3Y$ set instead
\[
 Z=hD\{H(n-LB)+K\}+H\{h(Y^2-e)-DB\}.
\]
In every case the complete polynomial equation is
\begin{equation}\label{eq:newcompact}
 P=Q(Z^2-M^2GV^2)+(M^2G-Z^2)=0.
\end{equation}
Every operation is addition, subtraction or multiplication; no division,
root, exponent variable or hidden quantified intermediate occurs.
For proof only let $z=Z/M$, $\kappa=K/H$, $\rho=e/D$,
$t=B/h$, $W=1+hL$, and let $\sigma=0$ in the cubic charts and
$\sigma=s/q$ in the partial chart. The exact rational identity is
\begin{equation}\label{eq:compactcarrier}
 z=n+\kappa+\sigma-Wt-\rho.
\end{equation}
In the partial chart $D\sigma=Y^2$ is the crucial reuse identity.

\begin{theorem}\label{thm:compact-six}
Each of the six polynomials above has a strictly positive integer zero
if and only if the retained syntactic calculus has a finite ZFC
refutation. Consequently, by the retained whole indexed calibration,
ZFC proves that its positive unsolvability is equivalent to
$\operatorname{Con}(\mathrm{ZFC})$.
\end{theorem}
\begin{proof}
We give the changed unbounded arithmetic argument, retaining the stated
universal Pell, ratio and native-calibration interfaces. Write
$\psi_A(r),\chi_A(r)$ for Pell ordinate and trace.

\emph{Roots before signs.}
Positive inputs imply $q,X,Y\ge32$, $c>k\ge2$, $A>4$, $U\ge c^2$,
$E>1$, $a>c$, $Q>c^2$, and $M>0$. The identities
\[
 F_{\rm sq}=(1+2a)^2,\qquad
 Q=4\Delta^2U^2G,\qquad \gcd(E,F_0)=1
\]
are unconditional. For any integral $b$, $V\ne0$ and $QV^2-1>0$.
Equation~\eqref{eq:newcompact} gives
$(Q-1)z^2=G(QV^2-1)$. Since $G\mid Q$, both other integer factors
are $-1$ modulo $G$; prime valuations force $G=g^2$.
Coprimality gives $E=\tau^2,F_0=f_0^2,g=\tau f_0$ with positive roots.
The rational number $y=|z|/g$ satisfies
\begin{equation}\label{eq:adjacent-new}
 QV^2-(Q-1)y^2=1.
\end{equation}
Rational adjacent-Pell descent makes $y$ integral: if its reduced
denominator is $m$, then $m^2\mid Q-1$; the norm-preserving map
\[
 (v,y)\longmapsto((2Q-1)v-2(Q-1)y, (2Q-1)y-2Qv)
\]
keeps the first coordinate integral, keeps the reduced denominator,
and strictly decreases $|v|>1$ after sign reflection. At $|v|=1$,
$y=1$. Thus $z=\varepsilon gy$ is integral.

Put $f=1+2a,I=2Ug$. Then
$f^2=F_{\rm sq}$, $f^2-\Delta I^2=1$, $c^2\mid I$ and
$Q=(\Delta I)^2$. The retained signed all-index folding theorem,
which allows either sign of $V$, applies to these equations and
$V=-c\pmod f$. It gives $c=\psi_A(p)$ for an odd $p\ge1$.
The first norm gives $k=2\psi_{2v+1}(n_0)$, $n_0\ge1$.
The value $n_0=1$ would give $c=2Y+1$ strictly between
$\psi_A(1)=1$ and $\psi_A(2)=2A$. Hence
\begin{equation}\label{eq:earlyc}
 k\ge8XY^2+4,\qquad c>kY>8XY^3\ge8q^4,
 \qquad \kappa>0.
\end{equation}
No sign of $B,C,R,b$ has been used.

\emph{Sign-free transport.}
At a zero, $b=0$ would make
$y^2=c^2+(c^2-1)/(Q-1)$ lie between consecutive integers. Thus
$\operatorname{sgn}V=\operatorname{sgn}b$ and $|V|>|b|$.
Also $a^2-G=E[\Delta U^2E(\Delta U^2-1)-1]>0$,
$F_{\rm sq}-c>3a>3g$, and $y<3|V|/2$.
Therefore $gy<V^2/2$.

For $L=bV$ a stronger bound is needed. Suppose $p=3$, so
$c=4A^2-1$. If $n_0\ge3$, then $kY>32X^2Y^5$, whereas
$c<16X^2Y^2$, a contradiction. Hence $n_0=2$ and $k=8XY^2+4$.
For the integer $u=X+2-2Y$ one obtains
\[
 \eta=4Y^2(Xu+1)+8Y(X+1)+3-4Y.
\]
If $u\le-1$ this is negative, since
$Y(X-1)>2(X+1)$ for $X,Y\ge32$.
If $u\ge2$ it exceeds $k$. Hence $u\in\{0,1\}$, or
$2Y-X\in\{1,2\}$, impossible since $q\mid X,Y$.
Thus $p\ge5$ and $c>c_3=4\Delta+3$.

Writing $S=\Delta I$, Pell classification of
\eqref{eq:adjacent-new} gives $|V|=\chi_S(j)/S$ for odd $j$.
The first possibility is $|V|=1$, already excluded. For the third,
$V=\pm(4Q-3)$ gives respectively
\[
 F_{\rm sq}(b-4\Delta)=c-c_3,\qquad
 F_{\rm sq}(b+4\Delta)=c+c_3.
\]
Both right sides lie strictly between $0$ and $F_{\rm sq}$.
Thus $j\ge5$, $|V|>8Q^2$, and, since $Q>F_{\rm sq}>2c$,
\[
 |b|>\frac{8Q^2-c}{F_{\rm sq}}>4F_{\rm sq}>4g.
\]
Consequently $L=bV=|b||V|>0$ and $gy<3L/8<L/2$.

For either weight $e<b$ because $b-e=cT-k+1>0$, so
$\rho<|b|/D$. If $B\le-1$, then $Wt<-L$; the carrier gives
\[
 n< L/2-L+|b|/D<0,
\]
using $D>2$, $|V|>|b|$ and $|V|>1$. Contradiction.
Thus $B\ge0$, and $F>0$ forces $C>0$.
Only now use the mask estimate $3q<R<q^4-2$.
Equation~\eqref{eq:earlyc} gives $c>2R$, hence $b>0,V>0$.
The retained normalized-rank and signed-index theorems, using
$c^2\mid I$, $c>2R,c>2p$ and the fixed-minus residues, give $p=R$.
Their hypotheses do not require the removed shift $T+\alpha$.
Odd $R$ then gives even $q$ and odd $J$.

\emph{Recover a small index before the full denominator.}
The recurrence at $P_0=2XY^2+1$ gives
\[
 k\equiv2n_0\pmod {4XY^2}.
\]
Since $P_0>A$, $n_0\ge R$ would imply $k/2\ge c$.
Thus $1\le n_0<R$. Integral $z$ gives $D\mid hHe$ in the cubic
charts, where $q^4\mid D$. Since $H=-1\pmod q$ and $h=q-1$,
this yields $q^4\mid e$.
In the partial chart clearing $qhH$ gives $D\mid qhHe$;
$q^5\mid D$ again gives $q^4\mid e$.
Crucially $q^4\mid4XY^2$ in \emph{all} these charts.
Therefore $q^4\mid2n_0-R-1$, and
$|2n_0-R-1|<R<q^4$ forces $2n_0=R+1$.
The full recurrence now gives $4XY^2\mid e$.
Each selected $D$ divides $4XY^2$; for $D=q^3Y$ the quotient is
$4ws$. Thus $\rho$ is integral. No false coprimality of $Y$ and $h$
was invoked.

\emph{Projection and source decoding.}
Write $R=2r+1$, so $r\ge48$. The retained ratio theorem at scale one,
with $\xi=(X+1)^{2r}/X^r$, gives
\[
 \xi<\frac{\psi_A(2r+1)}{\psi_{P_0}(r+1)}
 <\xi\left(1+\frac1{Y(X+1)}\right)^{2r}.
\]
Together with $Y<c/k<Y+1$, this gives $Y>X^r/3$ and
$X^{r+1}<H$. With $\rho$ integral the carrier gives
$H\mid hK$ in the cubic charts, $H\mid qhK$ in the partial chart.
Either implies the inherited premise $H\mid q^3hK$.
Universal projection and odd $H$ give
\[
 H\mid q^3h(X-2^R)(2^RX+1).
\]
The retained whole-product bound puts the absolute value of this
\emph{entire} product below $X^{r+1}<H$, hence $X=2^R$.
In particular $\kappa$ is integral; no individual nonunit factor was
cancelled modulo $H$.
The cubic carrier gives $h\mid B$. The partial carrier gives
$s/q-WB/h\in\Z$; coprimality of $q,h$ gives $q\mid s$ and $h\mid B$.
Thus $q^3\mid Y$ and $t\ge0$ is integral in every chart.
The retained fixed-minus parity and half-binomial rounding theorems
now have their full hypotheses. Their fresh reconstruction at
$A_{\rm source}=Y(X+1)+2$, followed by the whole indexed native
calibration, supplies a finite ZFC refutation.

\emph{Positive converse.}
A refutation supplies the retained canonical outer data,
$X=2^R$, $R\equiv3\pmod4$, $R>3q$, $q^3\mid Y$, integral $t\ge0$,
and the required half-binomial $Y$. Set $s=Y/q^3$ in cubic charts,
$s=Y/q^2$ in the partial chart. Form the scale-one bases and canonical
$c=\psi_A(R)$, $k=2\psi_{P_0}((R+1)/2)$, $\tau=\chi_{P_0}((R+1)/2)$.
The ratio/rounding theorem gives $0<\eta=c-kY<k$, so
$\zeta=k-\eta>0$. Also $4XY^2\mid e$, hence $\rho\ge0$ is integral,
and projection gives integral $0<\kappa<Xc$.
Choose a positive matrix-period multiple at $A$ modulo $c^2\tau$,
with seed pair $(f_0,I_0)$ and $c^2\tau\mid I_0$.
Put $i=I_0/(c^2\tau)>0$ and double the pair; this produces exactly
$f=2f_0^2-1$, $I=2f_0I_0$, and all literal root definitions above.
The retained auxiliary completeness theorem supplies positive $V,y$ and
\[
 T=\frac{V+c+Rf^2}{cf^2}\in\Z_{>0}.
\]
No subtraction of $\alpha$ is needed. Set
$n=gy+Wt+\rho-\kappa-\sigma$.
In the partial chart $\sigma=Y/q^3$ is now integral; in cubic charts
it is zero. Since $\sigma<Y<Xc$, $\kappa<Xc$, and
$f_0>(A-1)I_0>2Xc$, we have $g>\kappa+\sigma$ and therefore $n>0$.
The carrier is $z=gy$, proving~\eqref{eq:newcompact}. All independent
coordinates are positive; this is a fresh witness construction, not an
unjustified positive inverse shear.
\end{proof}

\subsection{Exact degree and reproducibility}
Let $q_0=\beta J$, $k_0=\eta+\zeta$, $C_0=q_0-F-\alpha$.
The full highest homogeneous forms are respectively
\begin{align*}
 (XY,V^2):&\quad2^{14}q_0^{378}w^{48}s^{110}i^{20}k_0^{64}T^4C_0^2,\\
 (XY,bV):&\quad-2^{10}q_0^{284}w^{36}s^{82}i^{14}k_0^{46}T^2,\\
 (q^4,V^2):&\quad2^{14}q_0^{378}w^{46}s^{108}i^{20}k_0^{64}T^4C_0^2,\\
 (q^4,bV):&\quad-2^{10}q_0^{284}w^{34}s^{80}i^{14}k_0^{46}T^2,\\
 (q^3Y,V^2):&\quad2^{14}q_0^{272}w^{46}s^{110}i^{20}k_0^{64}T^4C_0^2,\\
 (q^3Y,bV):&\quad-2^{10}q_0^{206}w^{34}s^{82}i^{14}k_0^{46}T^2.
\end{align*}
For $V^2$ the unique top summand is $QZ^2$; for $bV$ it is
$-QM^2GV^2$. The displayed nonzero forms prove the claimed exact global
degrees at the actual positive $\beta$. Every binary operation is charged,
including squares and multiplication by scalar literals. All inputs,
ports and gates are live. Six deterministic DAGs, frozen parent data,
coefficient bytes, recipe and mathematical provenance accompany the proof.
The independent checker covers every row by exact polynomial-cut identities,
propagates complete leading forms, authenticates all fixed bytes, and rejects
48 deliberate algebraic corruptions. Normal and optimized Python replays
and deterministic rebuilds are supplied. These finite checks are not
substitutes for the unbounded theorem above.

A tempting $Y=qs,D=q^6$ route at $81/410$ or $82/304$ was rejected:
$4XY^2$ is only visibly divisible by $q^3$, insufficient for its index
argument. It is not an accepted equation or delivered endpoint.
